\documentclass[12pt]{article}

\usepackage[margin=1in]{geometry}
\usepackage{enumitem}
\usepackage{amsmath}
\usepackage{amssymb}
\usepackage{array}
\usepackage{titlesec}

\title{Lecture 5 Notes: \\ Aristotle's \textit{Prior Analytics}, Book II}
\author{}
\date{}

\begin{document}

\maketitle

\section*{Central Theme}

Aristotle's \textit{Prior Analytics}, Book II studies what happens after the
basic machinery of syllogism has been constructed.

Book I explained how syllogisms are formed. Book II explains how syllogisms
behave, how they fail, how they are refuted, and how syllogistic reasoning is
related to induction, example, objection, and enthymeme.

\[
\boxed{
\textit{Prior Analytics}, \text{ Book II}
=
\text{the behavior, defects, and relatives of syllogism}
}
\]

\section*{1. The Transition from Book I to Book II}

Book I gave Aristotle's construction of syllogistic logic.

\[
\text{Book I: How syllogisms are formed}
\]

\[
\text{Book II: How syllogisms behave, fail, and extend into other arguments}
\]

Aristotle begins Book II by reminding us that the figures, premises, modes of
proof, and methods of finding arguments have already been explained. Book II
therefore does not begin again. It applies, tests, and extends what has already
been established.

\subsection*{Teaching Point}

Book II is Aristotle's applied theory of syllogism.

\[
\boxed{
\text{Book I builds the syllogism; Book II studies syllogism in action.}
}
\]

\section*{2. The Three Parts of Book II}

Book II has three broad movements.

\[
\begin{array}{ll}
\textbf{I. Properties} & \text{how syllogisms behave} \\
\textbf{II. Defects} & \text{how syllogisms fail or mislead} \\
\textbf{III. Arguments akin to syllogism} & \text{induction, example, objection, enthymeme}
\end{array}
\]

\subsection*{Teaching Point}

Book II is not merely a supplement. It shows that logic must account for proof,
error, debate, refutation, and persuasion.

\[
\boxed{
\text{A theory of reasoning must explain both valid proof and failed proof.}
}
\]

\section*{3. More Than One Conclusion from the Same Premises}

Aristotle begins with a property of syllogisms: many syllogisms yield more than
one conclusion.

Universal syllogisms produce more than one result. Particular affirmative
syllogisms also yield more than one result. Particular negative syllogisms yield
only the stated conclusion.

The reason is conversion.

\[
\text{All propositions convert except the particular negative.}
\]

\subsection*{Examples}

If we prove:

\[
A \text{ belongs to all } B,
\]

then by conversion we may also infer:

\[
B \text{ belongs to some } A.
\]

If we prove:

\[
A \text{ belongs to no } B,
\]

then we may also infer:

\[
B \text{ belongs to no } A.
\]

But if we prove:

\[
A \text{ does not belong to some } B,
\]

we cannot infer:

\[
B \text{ does not belong to some } A.
\]

\subsection*{Whiteboard}

\[
\begin{array}{ll}
\text{Universal syllogism} & \Rightarrow \text{more than one conclusion} \\
\text{Particular affirmative} & \Rightarrow \text{more than one conclusion} \\
\text{Particular negative} & \Rightarrow \text{only the stated conclusion}
\end{array}
\]

\subsection*{Teaching Point}

Conversion lets one proof generate additional conclusions.

\[
\boxed{
\text{The logical force of a syllogism may exceed its stated conclusion.}
}
\]

\section*{4. True Conclusions from False Premises}

Aristotle next examines the relation between the truth of premises and the truth
of conclusions.

His central claims are:

\[
\boxed{
\text{A false conclusion cannot follow from true premises.}
}
\]

\[
\boxed{
\text{A true conclusion can follow from false premises.}
}
\]

\subsection*{Why a False Conclusion Cannot Follow from True Premises}

If the premises are true and the syllogism is valid, then the conclusion follows
of necessity. If the conclusion were false, the same thing would both be and not
be in the same respect, which is impossible.

\subsection*{Why a True Conclusion Can Follow from False Premises}

False premises can accidentally lead to a true conclusion.

For example:

\[
\text{All stones are animals.}
\]

\[
\text{All humans are stones.}
\]

\[
\therefore \text{All humans are animals.}
\]

The conclusion is true, but both premises are false.

\subsection*{Teaching Point}

The truth of the conclusion is not the same as the soundness of the proof.

\[
\boxed{
\text{Truth of conclusion} \neq \text{truth of explanation}
}
\]

\[
\boxed{
\text{A syllogism may be formally valid while failing as demonstration.}
}
\]

\section*{5. Fact and Reason}

Aristotle says that a true conclusion drawn from false premises may be true
``in respect to the fact,'' but not ``in respect to the reason.''

\[
\text{Fact: the conclusion is true.}
\]

\[
\text{Reason: the premises explain why it is true.}
\]

\subsection*{Example}

\[
\text{All humans are animals.}
\]

This is true as a fact.

But if we try to prove it by saying:

\[
\text{All stones are animals.}
\]

\[
\text{All humans are stones.}
\]

then we have not given the reason why humans are animals.

\subsection*{Teaching Point}

Demonstration is not satisfied with true conclusions. It requires true
explanatory premises.

\[
\boxed{
\text{Scientific proof must show not only that something is so, but why it is so.}
}
\]

\section*{6. Circular or Reciprocal Proof}

Aristotle then discusses circular proof.

Circular proof means proving by means of the conclusion. One converts one of
the original premises and uses the earlier conclusion to infer a premise that
had originally been assumed.

\subsection*{Pattern}

Original proof:

\[
A \text{ belongs to } C \text{ through } B.
\]

Circular proof:

\[
\text{Use } A \text{ belongs to } C
\]

together with a converted premise to prove one of the original premises.

\subsection*{Teaching Point}

Circular proof is possible only when the terms are convertible.

\[
\boxed{
\text{Circular proof depends on convertibility.}
}
\]

\section*{7. Why Circular Proof Matters}

Aristotle is not merely saying that circular proof is always useless. He is
analyzing when reciprocal proof can occur.

But in strict demonstration, circularity remains a problem if it reverses the
proper order of explanation.

\[
\text{Demonstration should proceed from what is prior and better known.}
\]

If the conclusion is used to prove the premises, the explanatory order becomes
unclear.

\subsection*{Teaching Point}

Reciprocal proof may show logical convertibility, but it does not always give
scientific explanation.

\[
\boxed{
\text{Logical reciprocity is not always explanatory priority.}
}
\]

\section*{8. Conversion of Syllogisms}

To convert a syllogism is to alter the conclusion and then form a new syllogism
showing that one of the original premises must be destroyed.

\subsection*{Basic Pattern}

\[
\text{Premise 1} + \text{Premise 2} \Rightarrow \text{Conclusion}
\]

Then:

\[
\text{Opposite of conclusion} + \text{Premise 1}
\Rightarrow
\text{Denial of Premise 2}
\]

or:

\[
\text{Opposite of conclusion} + \text{Premise 2}
\Rightarrow
\text{Denial of Premise 1}
\]

\subsection*{Teaching Point}

Conversion shows how a conclusion can be used against one of the premises.

\[
\boxed{
\text{To convert a syllogism is to use its conclusion against its premises.}
}
\]

\section*{9. Contrary and Contradictory Conversion}

Aristotle distinguishes converting the conclusion into its contrary from
converting it into its contradictory.

\[
\text{Contrary opposition: all / none}
\]

\[
\text{Contradictory opposition: all / not all}
\]

\subsection*{Example}

Original:

\[
\text{Every human is animal.}
\]

Contrary:

\[
\text{No human is animal.}
\]

Contradictory:

\[
\text{Not every human is animal.}
\]

\subsection*{Teaching Point}

The kind of opposition matters. A contrary and a contradictory do not have the
same logical force.

\[
\boxed{
\text{Changing the form of opposition changes the refutation.}
}
\]

\section*{10. Reductio ad Impossibile}

A syllogism \textit{per impossibile} is an argument by reduction to
impossibility.

\subsection*{Basic Pattern}

\[
\text{Assume the contradictory of what is to be proved.}
\]

\[
\text{Add an accepted premise.}
\]

\[
\text{Derive an impossibility.}
\]

\[
\text{Reject the assumption.}
\]

\[
\text{Therefore, accept the original claim.}
\]

\subsection*{Whiteboard}

\[
\text{Assume not-}P
\]

\[
\text{Derive } \bot
\]

\[
\therefore P
\]

\subsection*{Teaching Point}

Reductio proves a claim indirectly by showing that its contradictory leads to
impossibility.

\[
\boxed{
\text{Reductio is indirect proof through contradiction.}
}
\]

\section*{11. Why Reductio Requires the Contradictory}

Aristotle insists that in a reductio one must assume the contradictory, not
merely the contrary.

Contradictories cannot both be true and cannot both be false. Therefore, if the
contradictory is refuted, the original claim must be true.

Contraries cannot both be true, but they may both be false. Therefore, if the
contrary is refuted, the original claim is not necessarily established.

\subsection*{Whiteboard}

\[
\text{Contradictories: exactly one must be true.}
\]

\[
\text{Contraries: both cannot be true, but both may be false.}
\]

\subsection*{Teaching Point}

\[
\boxed{
\text{Reductio works through contradiction, not mere contrariety.}
}
\]

\section*{12. Ostensive Proof and Proof Per Impossibile}

Aristotle compares direct proof with proof by impossibility.

\subsection*{Ostensive Proof}

Ostensive proof begins from accepted premises and directly establishes the
conclusion.

\[
P_1 + P_2 \Rightarrow C
\]

\subsection*{Proof Per Impossibile}

Proof \textit{per impossibile} begins by assuming the contradictory of the
desired conclusion and deriving an impossibility.

\[
\text{not-}C \Rightarrow \bot
\]

\[
\therefore C
\]

\subsection*{Teaching Point}

Aristotle argues that ostensive proof and proof by impossibility are closely
related. What can be proved directly can also be proved indirectly, and what can
be proved indirectly can be connected back to direct proof through the same
terms.

\[
\boxed{
\text{Direct and indirect proof are different routes through the same logical terrain.}
}
\]

\section*{13. Reasoning from Opposites}

Aristotle next examines whether conclusions can be drawn from opposed premises.

He distinguishes several kinds of opposition:

\[
\begin{array}{ll}
\text{Universal affirmative} & \text{Every science is good} \\
\text{Universal negative} & \text{No science is good} \\
\text{Particular affirmative} & \text{Some science is good} \\
\text{Particular negative} & \text{Some science is not good}
\end{array}
\]

\subsection*{Contraries}

\[
\text{Every science is good.}
\]

\[
\text{No science is good.}
\]

These are contraries.

\subsection*{Contradictories}

\[
\text{Every science is good.}
\]

\[
\text{Some science is not good.}
\]

These are contradictories.

\subsection*{Teaching Point}

Logical opposition is not one simple relation. Different kinds of opposition
produce different kinds of argumentative results.

\[
\boxed{
\text{Opposition in premises changes what sort of conclusion can be drawn.}
}
\]

\section*{14. Defects of Syllogism}

The second major part of Book II discusses defects of reasoning.

The three central defects are:

\[
\begin{array}{ll}
1. & \text{Begging the question} \\
2. & \text{False cause} \\
3. & \text{Error from false premises or faulty application}
\end{array}
\]

\subsection*{Teaching Point}

Aristotle does not only classify good arguments. He also diagnoses bad ones.

\[
\boxed{
\text{Logic must explain failure as well as success.}
}
\]

\section*{15. Begging the Question}

Begging the question occurs when someone tries to prove what is not self-evident
by means of itself.

\[
\boxed{
\text{Begging the question} =
\text{proving what is not self-evident by itself}
}
\]

This may happen directly, by assuming the very thing to be proved. It may also
happen indirectly, by proving a claim through something that naturally depends
on that claim.

\subsection*{Pattern}

\[
A \text{ is proved by } B
\]

\[
B \text{ is proved by } C
\]

\[
C \text{ actually depends on } A
\]

\[
\therefore A \text{ has been used to prove itself.}
\]

\subsection*{Teaching Point}

A proof fails if its premises depend on its conclusion.

\[
\boxed{
\text{A conclusion cannot be the hidden source of its own proof.}
}
\]

\section*{16. Begging the Question and Demonstration}

Begging the question violates the order of demonstration.

Demonstration should proceed from what is:

\begin{itemize}[itemsep=0.25em]
    \item prior,
    \item better known,
    \item more certain,
    \item explanatory.
\end{itemize}

But begging the question proceeds from what is equally uncertain or dependent
on the conclusion.

\subsection*{Teaching Point}

Begging the question is not merely a debating trick. It is a failure of
scientific order.

\[
\boxed{
\text{Demonstration requires explanatory priority.}
}
\]

\section*{17. False Cause}

The objection ``this is not the reason why'' belongs especially to reductio
arguments.

False cause occurs when an impossible conclusion is said to follow from a
hypothesis, but the impossibility would have followed even without that
hypothesis.

\subsection*{Pattern}

\[
\text{Assume } H
\]

\[
\text{Derive impossibility } \bot
\]

But if:

\[
\bot \text{ follows even without } H,
\]

then:

\[
H \text{ was not the cause of the impossibility.}
\]

\subsection*{Teaching Point}

In a reductio, the contradiction must depend on the assumption being refuted.

\[
\boxed{
\text{Not every contradiction proves what we want it to prove.}
}
\]

\section*{18. Falsity in Premises}

Aristotle argues that a false argument depends on the first false statement in
it.

If a false conclusion is drawn, then one or more of the premises must be false.

\[
\boxed{
\text{False conclusion} \Rightarrow \text{at least one false premise}
}
\]

This follows because a false conclusion cannot be drawn from true premises in a
valid syllogism.

\subsection*{Teaching Point}

Error in a syllogism must be traced to a false premise or to an invalid
structure.

\[
\boxed{
\text{To diagnose a false conclusion, trace the argument back to its first false premise.}
}
\]

\section*{19. Impeding Opposing Arguments}

Aristotle includes tactical advice for argument.

If one wants to avoid having a syllogism drawn against oneself, one should
avoid granting the same term twice. The repeated term becomes the middle term,
and the middle term allows the opponent to construct a syllogism.

\subsection*{Whiteboard}

\[
\text{Repeated term} = \text{middle term}
\]

\[
\text{Middle term} = \text{danger point in debate}
\]

\subsection*{Teaching Point}

Dialectical strategy depends on controlling the middle term.

\[
\boxed{
\text{If you grant the middle, you may grant the argument.}
}
\]

\section*{20. Concealing One's Own Argument}

Aristotle also explains how an attacker may conceal an argument.

Instead of asking for connected premises in obvious order, one should ask for
premises in a less obvious sequence, leaving the conclusion hidden until the
middle terms are secured.

\subsection*{Pattern}

Obvious order:

\[
A \rightarrow B \rightarrow C \rightarrow D
\]

Concealed order:

\[
A \rightarrow B,
\quad
D \rightarrow E,
\quad
B \rightarrow C,
\quad
C \rightarrow D.
\]

\subsection*{Teaching Point}

In live dialectic, the order of questioning can hide the structure of the
syllogism.

\[
\boxed{
\text{Argument is not only form; it is also strategy.}
}
\]

\section*{21. Refutation}

A refutation is a syllogism that establishes the contradictory of the opponent's
claim.

\[
\boxed{
\text{Refutation} = \text{a syllogism proving the contradictory}
}
\]

Refutation is possible only when enough premises are granted to construct a
syllogism.

\[
\text{No concession} \Rightarrow \text{no syllogism}
\]

\[
\text{No syllogism} \Rightarrow \text{no refutation}
\]

\subsection*{Teaching Point}

Refutation requires premises from the opponent.

\[
\boxed{
\text{You cannot refute someone syllogistically unless they grant material for a syllogism.}
}
\]

\section*{22. Error and Knowledge}

Aristotle discusses how error can arise in thought.

Sometimes a person has universal knowledge but fails to apply it correctly to a
particular case.

\subsection*{Example}

A person may know:

\[
\text{Every triangle has angles equal to two right angles.}
\]

But if the person does not recognize this figure as a triangle, then he does
not actively know that this figure has angles equal to two right angles.

\subsection*{Teaching Point}

Universal knowledge is not the same as actual application to a particular case.

\[
\boxed{
\text{Error often arises not from lacking a universal, but from failing to apply it.}
}
\]

\section*{23. Knowledge as Possession and Exercise}

Aristotle distinguishes different senses of knowing.

One may:

\begin{itemize}[itemsep=0.25em]
    \item possess universal knowledge,
    \item possess knowledge relevant to a particular,
    \item actively exercise that knowledge.
\end{itemize}

\subsection*{Teaching Point}

A person may know in one sense and fail to know in another.

\[
\boxed{
\text{Knowledge possessed is not always knowledge exercised.}
}
\]

\section*{24. Arguments Akin to Syllogism}

The final major part of Book II discusses forms of reasoning related to
syllogism.

These include:

\[
\begin{array}{ll}
1. & \text{Induction} \\
2. & \text{Example} \\
3. & \text{Reduction} \\
4. & \text{Objection} \\
5. & \text{Enthymeme}
\end{array}
\]

\subsection*{Teaching Point}

Aristotle expands from strict syllogistic demonstration to nearby forms of
reasoning used in inquiry, debate, and persuasion.

\[
\boxed{
\text{Syllogism is the center, but not the whole of reasoning.}
}
\]

\section*{25. Induction}

Aristotle says that every belief comes either through syllogism or from
induction.

Induction establishes a relation between one extreme and the middle by means of
the other extreme.

\subsection*{Syllogism}

\[
\text{Major term belongs to minor through the middle.}
\]

\subsection*{Induction}

\[
\text{Major term belongs to middle through particulars.}
\]

\subsection*{Example}

Let:

\[
A = \text{long-lived}
\]

\[
B = \text{bileless}
\]

\[
C = \text{particular long-lived animals: man, horse, mule}
\]

If all the particular bileless animals are long-lived, and the enumeration is
complete, then induction establishes that bileless animals are long-lived.

\subsection*{Teaching Point}

Induction moves through particulars.

\[
\boxed{
\text{Syllogism proceeds through the middle; induction proceeds through particulars.}
}
\]

\section*{26. Induction and What Is Clearer to Us}

Aristotle says that syllogism through the middle term is prior and better known
by nature, but induction is clearer to us.

\[
\text{Syllogism} = \text{prior by nature}
\]

\[
\text{Induction} = \text{clearer to us}
\]

\subsection*{Teaching Point}

Human learning often begins with particular cases, even if scientific
understanding ultimately seeks universal causes.

\[
\boxed{
\text{We often learn from particulars before we understand through causes.}
}
\]

\section*{27. Example}

An example is reasoning from one particular case to another similar particular
case.

\[
\boxed{
\text{Example} = \text{particular} \rightarrow \text{similar particular}
}
\]

\subsection*{Aristotle's Pattern}

Suppose we want to show that fighting the Thebans is bad for the Athenians.

We use the case of the Thebans fighting the Phocians, which was bad for the
Thebans.

Both cases fall under the same general description:

\[
\text{making war against neighbors}
\]

\subsection*{Structure}

\[
D \sim C \quad \text{under } B
\]

\[
A \text{ belongs to } D
\]

\[
\therefore A \text{ belongs to } C
\]

\subsection*{Teaching Point}

Example is not reasoning from whole to part or part to whole, but from part to
part under a shared type.

\[
\boxed{
\text{Example reasons from a known case to a similar case.}
}
\]

\section*{28. Reduction}

Reduction is an argument that moves us closer to knowledge without yet giving
complete demonstration.

It occurs when the major term clearly belongs to the middle, but the relation
of the middle to the minor is uncertain, though equally or more probable than
the conclusion.

\subsection*{Pattern}

\[
A \text{ clearly belongs to } B
\]

\[
B \text{ uncertainly belongs to } C
\]

\[
\therefore A \text{ may belong to } C
\]

\subsection*{Example}

\[
A = \text{what can be taught}
\]

\[
B = \text{knowledge}
\]

\[
C = \text{virtue}
\]

It is clear that knowledge can be taught. The uncertain question is whether
virtue is knowledge. If virtue is knowledge, then virtue can be taught.

\subsection*{Teaching Point}

Reduction does not complete knowledge. It brings the inquiry nearer to
knowledge.

\[
\boxed{
\text{Reduction is progress toward proof, not proof itself.}
}
\]

\section*{29. Objection}

An objection is a premise contrary to a premise.

\[
\boxed{
\text{Objection} = \text{a premise contrary to a premise}
}
\]

Aristotle distinguishes universal and particular objections.

\[
\begin{array}{ll}
\textbf{Universal objection} & \text{proved through the first figure} \\
\textbf{Particular objection} & \text{proved through the third figure}
\end{array}
\]

\subsection*{Example}

If someone says:

\[
\text{Contraries are subjects of a single science,}
\]

one may object universally:

\[
\text{Opposites are never subjects of a single science.}
\]

Or one may object particularly:

\[
\text{The knowable and the unknowable are not subjects of a single science.}
\]

\subsection*{Teaching Point}

An objection attacks a premise, not merely the conclusion.

\[
\boxed{
\text{To object well, strike the premise that supports the proof.}
}
\]

\section*{30. Enthymeme}

Book II ends with the enthymeme.

Aristotle connects enthymeme with probability and sign.

\subsection*{Probability}

A probability is a generally approved proposition concerning what happens for
the most part.

\[
\text{The envious hate.}
\]

\[
\text{The beloved show affection.}
\]

\subsection*{Sign}

A sign is an indicator: something such that when it is, another thing is, or
when it has come to be, another thing has come to be before or after.

\[
\text{Having milk may be a sign of being with child.}
\]

\subsection*{Enthymeme}

An enthymeme is a syllogism from probabilities or signs.

\[
\boxed{
\text{enthymeme} = \text{syllogism from probabilities or signs}
}
\]

\subsection*{Teaching Point}

The enthymeme is syllogistic reasoning adapted to persuasion.

\[
\boxed{
\text{Rhetorical reasoning does not abandon logic; it uses probable logic.}
}
\]

\section*{31. Signs in the Figures}

Aristotle says signs may be taken according to the position of the middle term
in the figures.

\subsection*{First Figure Sign}

Example:

\[
\text{A woman has milk.}
\]

\[
\text{Therefore she is with child.}
\]

Here the sign functions as a middle term.

\subsection*{Third Figure Sign}

Example:

\[
\text{Pittacus is good.}
\]

\[
\text{Pittacus is wise.}
\]

\[
\text{Therefore wise men are good.}
\]

This is weaker because one particular wise man does not prove something about
all wise men.

\subsection*{Middle Figure Sign}

Example:

\[
\text{Women with child are pale.}
\]

\[
\text{This woman is pale.}
\]

\[
\text{Therefore this woman is with child.}
\]

This is refutable because paleness may have another cause.

\subsection*{Teaching Point}

Signs may persuade, but not all signs prove equally well.

\[
\boxed{
\text{A sign is strongest when it functions through the right middle term.}
}
\]

\section*{32. Final Architecture of Book II}

Book II moves from the internal behavior of syllogisms to the broader world of
argument.

\[
\text{properties}
\rightarrow
\text{truth and falsity}
\rightarrow
\text{circular proof}
\rightarrow
\text{conversion and reductio}
\rightarrow
\text{defects}
\rightarrow
\text{refutation}
\rightarrow
\text{arguments akin to syllogism}
\]

\subsection*{Teaching Point}

Book II shows syllogism at work in proof, debate, error, discovery, and
persuasion.

\[
\boxed{
\text{Book II is Aristotle's study of syllogism under pressure.}
}
\]

\section*{33. Summary of Main Ideas}

\begin{enumerate}[itemsep=0.5em]
    \item Many syllogisms produce more than one conclusion because of
    conversion.

    \item A false conclusion cannot follow from true premises, but a true
    conclusion can follow from false premises.

    \item A true conclusion from false premises may be true as a fact, but not
    as an explanation.

    \item Circular proof depends on convertibility, but may fail to preserve
    explanatory priority.

    \item Conversion of syllogisms uses an altered conclusion to refute one of
    the original premises.

    \item Reductio ad impossibile proves indirectly by assuming the
    contradictory and deriving an impossibility.

    \item Reductio requires contradiction, not mere contrariety.

    \item Begging the question occurs when what is to be proved is used,
    directly or indirectly, as its own proof.

    \item False cause occurs when the contradiction in a reductio does not
    depend on the assumption being refuted.

    \item Error often arises when universal knowledge is not applied to a
    particular case.

    \item Refutation is a syllogism proving the contradictory of an opponent's
    claim.

    \item Induction, example, reduction, objection, and enthymeme are arguments
    related to syllogism but distinct from strict demonstration.
\end{enumerate}

\section*{Final Framing}

\[
\boxed{
\textit{Prior Analytics}, \text{ Book II, shows syllogism in action.}
}
\]

Book I builds the syllogism. Book II tests it, attacks it, diagnoses its
failures, and shows how it extends into induction, example, objection, and
rhetorical persuasion.

\[
\boxed{
\text{A complete logic must explain proof, error, refutation, and persuasion.}
}
\]

\end{document}